\section{A complete real-zero classification}
Let $K=\Q(\sqrt{-47})$ and let $H$ be its Hilbert class field. The five reduced ideal classes occur in inverse pairs and the principal class. Representatives, up to inversion, are the forms
\[
 Q_0=x^2+xy+12y^2,\qquad Q_1=2x^2+xy+6y^2,\qquad
 Q_2=3x^2+xy+4y^2 .
\]
Reduction gives $(1,1,12),(2,\pm1,6),(3,\pm1,4)$. Thus the class group is cyclic of order five. Complex conjugation acts by inversion, and $\operatorname{Gal}(H/\Q)$ is dihedral of order ten. Put $B_i(z)=\sum_{x,y\in\Z} q^{Q_i(x,y)}$ and $q=e^{2\pi iz}$. Let $\lambda=(1+\sqrt5)/2$ and $\mu=(1-\sqrt5)/2$. The two normalized eigenforms and their difference are
\begin{align}
 f_1&=(B_0-\mu B_1-\lambda B_2)/2, &
 f_2&=(B_0-\lambda B_1-\mu B_2)/2,\nonumber\\
 E(z)&=\eta(z)\eta(47z)=(B_1-B_2)/2=(f_1-f_2)/\sqrt5. \label{eq:eta}
\end{align}
The $L_j$ of these forms are the Hecke $L$-functions of the two inverse pairs of nontrivial ideal class characters, or equivalently the two irreducible nonabelian Artin factors. They have the same gamma factor. See Berkovich and Patane, \href{https://arxiv.org/pdf/1302.2359}{Section 3, equations (3.1)--(3.4)}, and Buzzard, \href{https://www.ma.imperial.ac.uk/~buzzard/maths/research/papers/wt1.pdf}{Section 2.1}. We normalize
\[
 \Lambda_j(s)=\left(\frac{\sqrt{47}}{2\pi}\right)^s\Gamma(s)L_j(s),
 \qquad F_{c}(s)=\Lambda_1(s)-c\Lambda_2(s) \quad(c\in\R).
\]
All these completions are entire. A nonzero multiple of $\Lambda_2$, the one pencil direction omitted by this notation, will have no real zeros.

\begin{theorem}\label{th:real}
For $m\ge0$ let $C_m=\Lambda_1^{(2m)}(1/2)/\Lambda_2^{(2m)}(1/2)$. The thresholds strictly decrease to $1$. The ratio
\[
 R_m(t)=\frac{\Lambda_1^{(2m)}(1/2+t)}
                  {\Lambda_2^{(2m)}(1/2+t)}
\]
is positive and even. It decreases strictly from $C_m$ to $1$ as $t$ increases from zero to infinity, and $R_m''(0)<0$. Consequently the complete list of real zeros of $F_c^{(2m)}$ is as follows.
\begin{center}\begin{tabular}{ll}\toprule
Parameter & Real zeros with multiplicities \\\midrule
$c\le1$ or $c>C_m$ & None\\
$c=C_m$ & One double zero at $s=1/2$\\
$1<c<C_m$ & Two simple zeros, $s=1/2\pm t_c$, $t_c>0$\\\bottomrule
\end{tabular}\end{center}
For all real $s$, $\Lambda_1^{(2m)}(s)>\Lambda_2^{(2m)}(s)>0$.
\end{theorem}

\paragraph{1. Positivity of the smaller kernel.}
For $y\ge1/\sqrt{47}$, $q=e^{-2\pi y}<r=41/100$. The first terms of $f_2$ are $q-\lambda q^2-\mu q^3+\lambda q^4$. The absolute value of the $n$th coefficient is at most the number of ideals of norm $n$ in $K$, hence at most $d(n)\le n$, by the local split, inert and ramified cases. Dropping the positive third and fourth terms gives
\[
 \frac{f_2(iy)}q\ \ge\ 1-\tfrac53r-
 \frac{r^4(5-4r)}{(1-r)^2}>0.04.
\]
The eta product is positive on the imaginary axis, so $f_1(iy)>f_2(iy)$. Poisson summation on each binary lattice gives $B_i(i/(47y))=\sqrt{47}\,y B_i(iy)$: its Gram matrix has determinant $47$ and its adjugate is integrally equivalent to that matrix by rotation. Thus positivity and ordering extend to every $y>0$. This step gives nonvanishing at every real point for both completions by Mellin integration.

\paragraph{2. A modular function with just one pole.}
Define $h(z)=f_2(z)/E(z)$. The weights and characters cancel.
The Fricke transformations of the numerator and denominator agree, by the transformation just proved and \eqref{eq:eta} (then analytic continuation).
Hence $h$ descends to the compact curve $X_0(47)^+=X_0(47)/\langle W_{47}\rangle$. The eta function has no zeros in the upper half-plane. At infinity the numerator begins with $q$ and the denominator with $q^2$. As $47$ is prime, $X_0(47)$ has two cusps, and Fricke exchanges them. Their common image is unramified under this quotient and is the only pole of $h$, of order one. A meromorphic function with one simple pole gives a degree-one isomorphism with the sphere. This proves the coordinate property directly.

There are no elliptic identifications fixing $iy$ for $y>1/\sqrt{47}$. For a determinant-one matrix fixing $iy$ its diagonal entries are equal, and any nonidentity elliptic element has absolute trace less than two (working modulo sign). In $\Gamma_0(47)$ this would force both diagonals zero, inconsistent with the lower-left entry divisible by $47$ and determinant one. In the other coset write the matrix as
\[
 \begin{pmatrix}a&b\\47d&e\end{pmatrix}
 \begin{pmatrix}0&-1/\sqrt{47}\\\sqrt{47}&0\end{pmatrix}.
\]
Equality of diagonals gives $b=-d$, and the trace bound then forces $b=d=0$. The only resulting fixed point is $i/\sqrt{47}$. The quotient map is therefore locally unramified along the open ray. Since $h$ is a degree-one coordinate, $h(iy)$ has nonzero derivative there. It is real and tends to $+\infty$ like $q^{-1}$, so it increases strictly. In particular
\[
 b(u)=\frac{f_1(ie^u/\sqrt{47})}{f_2(ie^u/\sqrt{47})}
 =1+\frac{\sqrt5}{h(ie^u/\sqrt{47})}
\]
is strictly decreasing on $u>0$ and tends to one.

\paragraph{3. The ratio under the Mellin integral.}
Let $K_j(u)=e^{u/2} f_j(ie^u/\sqrt{47})$. The Fricke identities make these positive functions even. Their double exponential decay justifies all differentiations in
\[
 \Lambda_j(1/2+t)=\int_{\R}K_j(u)e^{tu}\,\dd u.
\]
Thus $R_m(t)$ is the expectation of $b(u)$ under the probability density proportional to $u^{2m}K_2(u)\cosh(tu)$ on $u>0$. Its derivative equals the covariance of $b(u)$ and $u\tanh(tu)$. For $t>0$ the second function increases strictly, while $b$ decreases strictly. Integrating $(b(u)-b(v))(u\tanh(tu)-v\tanh(tv))/2$ over two independent draws proves the covariance is strictly negative. At zero the same calculation after a second differentiation gives $R_m''(0)=\operatorname{Cov}(b(u),u^2)<0$. Weighting the measure at zero by $u^2$ proves $C_{m+1}<C_m$.

As $t\to\infty$ at fixed $m$, or $m\to\infty$ at $t=0$, these probability measures escape every compact interval. To check this, compare the mass on $[0,A]$ to that on $[A+1,A+2]$. The ratio tends to zero by the exponential tilt $\cosh(tu)$ or the power tilt $u^{2m}$, respectively. The bounded continuous $b$ tends to one there, so both limits are one. Since the denominator is strictly positive on the real line, strict derivative signs off the center and the strict second derivative at the center prove all the asserted multiplicities.\hfill$\square$

\paragraph{Odd derivatives.}
The real zeros of every odd derivative can be classified as well. For $t>0$, the ratio of derivatives of order $2m+1$ is the expectation of $b(u)$ under the density proportional to $u^{2m+1}K_2(u)\sinh(tu)$. Its derivative is the covariance with $u\coth(tu)$, a strictly increasing function of $u$ because $\cosh v\sinh v>v$ for $v>0$. The ratio decreases strictly from $C_{m+1}$ to $1$, by dominated convergence at zero and the same escape argument at infinity. The central zero of $F_c^{(2m+1)}$ is simple unless $c=C_{m+1}$, when it is exactly triple: its cubic Taylor coefficient is then nonzero because $C_{m+2}<C_{m+1}$. There is one further simple symmetric pair exactly for $1<c<C_{m+1}$. This completes the real-zero classification for all derivative orders.

\paragraph{A limitation illustrated by the same example.}
Taking $c=1$ gives an everywhere positive real completion for $E$. Nevertheless this particular completion has zeros with $\Re s>1$. Indeed $E$ is a combination of distinct Hecke eigenforms with different $T_2$ eigenvalues. Booker and Thorne's theorem for weight-one cusp forms therefore applies. Each such zero recurs in disjoint translates of a small disk by almost periodicity of the absolutely convergent Dirichlet series and Rouch\'e's theorem. The strictly positive Mellin kernel coexists with infinitely many nonreal zeros beyond the critical strip. These belong to the difference of the factors. See \href{https://doi.org/10.2140/ant.2014.8.2027}{Booker and Thorne, Theorem 1.1}.
